\begin{claim}{B.11f}{Proved}
Let $R$ be a henselian DVR with fraction field $K$ of characteristic
zero and residue field of characteristic $p>0$, with no perfection
assumption. Let $C/K$ be a smooth proper geometrically connected curve
of genus $h\ge2$, and fix one prime $\ell\ne p$. Define
$\mathcal P_\ell(C/R)$ by the following finite collection of conditions:

\begin{enumerate}
\item $\Jac(C)[p^\infty]$ extends to a $p$-divisible group over $R$.
\item For each connected geometric cyclic etale cover
$D\to C_{\bar K}$ of degree $\ell$, there is a finite extension
$L/K$ and a descent $D_L\to C_L$ such that
$\Jac(D_L)[p^\infty]$ extends to a $p$-divisible group over $R_L$.
\end{enumerate}
Here $R_L$ is the valuation ring of the unique extension of the
valuation. The extensions in the second condition may be ramified.
There are exactly $(\ell^{2h}-1)/(\ell-1)$ isomorphism classes of
these unlabelled geometric covers, so one common finite extension
can be used for all of them. Then
\[
          C\text{ has good reduction over }R
          \quad\Longleftrightarrow\quad\mathcal P_\ell(C/R).
\]
The conditions are conditions on integral $p$-adic Galois data of
$\Pi_C$ and its cyclic index-$\ell$ subgroups: their abelianized
pro-$p$ groups are the corresponding Jacobian Tate modules.
\end{claim}
\begin{proof}
In characteristic zero a $p$-divisible group over a field is an etale
system of finite group schemes, determined by its integral Tate module.
For the Jacobian this module is
$H_1(\Pi_C,\Z_p)$. Condition~(1) therefore asks for an integral
Barsotti--Tate extension of this representation. Condition~(2) asks
for potentially such extensions for a finite collection of open
geometric subgroups, with the arithmetic actions obtained after
descent of the covers. SGA~7~I, IX, Corollary~5.10 and
Proposition~5.13(a)--(b) equate these conditions respectively with good
and potentially good reduction of the Jacobians; at the residue prime
their hypothesis is $\operatorname{char}K=0$, without a perfect-residue
assumption \citep{BC_SGA7}. Different finite-field descents of a fixed
geometric cover become isomorphic over a further finite extension.
Thus the potential condition is independent of the descent chosen.

Connected cyclic degree-$\ell$ covers correspond to nonzero continuous
characters $\Pi_C\to\F_\ell$, up to multiplication by
$\F_\ell^*$. Since $H^1(C_{\bar K},\F_\ell)$ has dimension $2h$,
this gives the asserted finite count. Finite presentation descends
each cover and its deck action to a finite extension of $K$; taking
the compositum of the finitely many extensions proves the common-field
assertion. Inner automorphisms act trivially on their abelianizations,
so the resulting Tate-module actions do not depend on a point section.

We may prove the geometric assertion after completing $R$.
The natural $G_{\widehat K}\to G_K$ is an isomorphism by
\citet[Lemma 3.1]{BC_MinTsuj2022}, and it identifies inertia and
finite extensions. Characteristic-zero geometric comparison then
identifies the arithmetic fundamental groups, their finite etale
covers and the representations just described. Good reduction of a
Jacobian, and potential good reduction after a finite extension, are
preserved in both directions by this comparison: use the same inertia
action on a Tate module at an auxiliary prime different from $p$ and
the SGA~7 criterion. Good reduction of $C$ is likewise preserved in
both directions by Tamagawa's pro-$r$ outer-inertia criterion for
$r\ne p$ \citep[Theorem 5.3]{BC_Tamagawa}. Thus both sides of the
claimed equivalence are unchanged. This reduction ensures that the
proper lifting theorem SGA~1, X, Theorem~2.1 is used with its literal
\emph{complete local base} hypothesis.

If $C$ is good, its Jacobian is an abelian scheme, proving~(1).
For~(2), prime-to-$p$ specialization for a smooth proper curve
identifies each geometric degree-$\ell$ cover with a cover of the
geometric special fibre. That cover descends over a finite separable
residue extension and lifts over the corresponding finite unramified
extension of the complete trait. The exact inputs are SGA~1, X,
Corollary~3.9 and Theorem~2.1 \citep{BC_SGA1}.
Its lifted proper etale cover is smooth, so its Jacobian is good.
This proves the required potential condition for every geometric cover.

Conversely assume~(1) and~(2). By~(1) the Jacobian is good over $R$.
Take a finite extension over which $C$ has a stable model; stable
reduction applies over an arbitrary residue field, by properness of
the stable-curve stack \citep[Theorem 2.7]{BC_Knudsen}.
Pass to a further finite unramified extension when splitting geometric
components and nodes is necessary. Work on the geometric special fibre
$C_0$ and its connected dual graph $\Gamma$.
The toric rank of the Jacobian's Neron identity component is
$b_1(\Gamma)$. Precisely, SGA~7~I, IX, (12.3.6)--(12.3.7) and
(12.4.1) identify its character lattice with the kernel of the
oriented edge-to-vertex incidence map \citep{BC_SGA7}.
If the total stable model has a singularity $xy=\varpi^m$, its
regular resolution only subdivides that edge, leaving this kernel
unchanged; the same Picard--Neron comparison applies. Since the
Jacobian remains good under this base change, $b_1(\Gamma)=0$.
Thus $\Gamma$ is a tree.

Suppose there is a node $e$. Removing $e$ partitions $\Gamma$ into
two nonempty subtrees. Both sides contain a component of positive
genus. Indeed, if one side had $r$ rational components, its total
number of incident branches, including the one attachment at $e$,
would be $2(r-1)+1=2r-1$. Stability of every unmarked rational
component would require at least $3r$, a contradiction.
Choose one positive-genus component on each side and a nonzero
$\F_\ell$ character of its geometric fundamental group.

There is an etale $\F_\ell$-torsor $D_0\to C_0$ with these
characters on the two chosen components and zero characters on the
others. To verify this gluing assertion, normalization gives the exact
sequence of etale constant sheaves
\[
 0\longrightarrow\F_{\ell,C_0}\longrightarrow
 \nu_*\F_{\ell,\widetilde C_0}\longrightarrow
 \bigoplus_e\F_{\ell,e}\longrightarrow0.
\]
At a node the last map is the difference of the two branch values.
The skyscraper terms have no positive cohomology over the algebraically
closed residue field. The associated sequence therefore surjects
$H^1(C_0,\F_\ell)$ onto
$\bigoplus_v H^1(C_v,\F_\ell)$, with kernel
$H^1(\Gamma,\F_\ell)=0$. This constructs the desired class.
Its restriction to each side of $e$ is nonzero, so its cover of that
connected side is connected: a torsor under a group of prime order
is either connected or trivial. Above $e$ there are exactly $\ell$
nodes. Contracting the two connected lifted side graphs leaves two
vertices joined by $\ell$ edges. Consequently
\[
                     b_1(\Gamma_{D_0})\ge\ell-1>0.
\tag{B.pdiv.1}
\]

All the data of this finite etale torsor descend to a finite
\emph{separable} extension of the residue field: SGA~1, IX,
Theorem~4.10 makes finite etale covers invariant under a purely
inseparable extension, and finite presentation then descends them
from the separable closure \citep{BC_SGA1}. After the corresponding
unramified extension of the already chosen complete stable trait,
SGA~1, X, Theorem~2.1 lifts the cover and its $\F_\ell$-torsor
structure uniquely over the proper stable model. Its generic fibre
is geometrically connected: after any finite extension a disconnected
generic cover would disconnect the normal finite etale total cover,
whereas its connected geometric special fibre forbids that.
The lifted model is again stable, since an etale cover preserves
the nodal local equations and ampleness of the relative dualizing
bundle. It has the graph in (B.pdiv.1).

The geometric generic cover just obtained is one of those in~(2).
But its Jacobian has positive toric rank, and this remains true after
every finite field extension. Indeed the base-changed stable model
has the same geometric graph; resolving thicker nodes only subdivides
edges. The character-lattice formula therefore retains positive rank.
It cannot become an abelian scheme. A descent of this same geometric
cover satisfying~(2) would become isomorphic to it over a common finite
extension, contradicting this persistent positive toric rank. Thus
$C_0$ has no nodes, and $C$ has potentially good reduction.

It remains to descend good reduction, rather than merely assert
potential goodness. Choose a finite Galois stable extension and its
now smooth model. The finite stable-curve Isom scheme extends the
semilinear descent automorphisms, and specialization identifies their
actions on the prime-to-$p$ Tate module of the Jacobian
\citep[Theorem 1.11]{BC_DM}. By~(1), the original inertia acts
trivially on that Tate module. Each resulting special-curve
automorphism is therefore the identity. One can use
\citet[Theorem 1.13, Lemma 1.14]{BC_DM} directly: for a nonidentity
automorphism of a smooth curve of genus $h\ge2$, the graph has
nonnegative finite intersection with the diagonal, whereas trivial
action on $H^1$ would give intersection number $2-2h<0$ by their
Lefschetz--Weil calculation. This argument also covers wild finite
automorphisms.
Smooth specialization is functorial for these model automorphisms.
It follows that the entire original inertia acts trivially in the
outer group of the geometric pro-$r$ fundamental group, $r\ne p$.
Tamagawa's criterion now gives good reduction over the original trait.
This avoids an unsupported descent through a wild quotient.
\end{proof}

For the requested example $y^2=(x^3-5^6)(x^3-1)$ over $\Q_5$,
B.11 gives two elliptic special components and a separating node.
The Jacobian's $5$-divisible group extends, but a nontrivial cyclic
etale double cover on each elliptic component gives, by the preceding
construction, a geometric double cover of the genus-two generic curve
whose stable graph has two vertices and two edges. Its genus is three
and its Jacobian has toric rank one after every finite extension.
Thus condition~(2) of $\mathcal P_2$ fails, explicitly distinguishing
the curve from its good Jacobian. The graph and genus calculation is
included in \texttt{BC\_pdiv\_cover\_checks.py}.

\begin{claim}{B.12}{Proved}
Let $R$ be any henselian mixed-characteristic DVR, with residue
characteristic $p>0$, and let $U=\bar U\setminus D$ be hyperbolic.
Fix a prime $\ell\ge5$, $\ell\ne p$. Good marked reduction is
equivalent to the following criterion in $p$-adic data attached to
finite covers of $\Pi_U$:

There exist a finite \emph{unramified} extension $L/K$ and a
geometrically connected finite Galois cover $V\to U_L$ of degree
prime to $p$ whose smooth proper compactification $C$ has genus at
least two, whose extension $C\to\bar U_L$ is ramified at every
geometric cusp, and for which $\mathcal P_\ell(C/R_L)$ holds.

When $D$ is empty the identity cover is permitted. In the forward
direction the cover can have degree a power of $\ell$. No
perfection, completeness, semistability or chosen point hypothesis
is required. For perfect residue fields the crystalline finite-cover
criterion of B.11d--B.11e remains an alternative. A crystalline
comparison for the original open curve is not asserted.
\end{claim}
\begin{proof}
The proper group $\Pi_C$ is the quotient of the geometric subgroup
of the cover by the closed normal subgroup generated by its cusp
inertias, exactly as in (B.11d.1). Its cyclic index-$\ell$ subgroups
and the integral $p$-adic abelianizations used in B.11f are therefore
data of finite covers and quotients of $\Pi_U$, with their inherited
arithmetic actions. Finite extensions and good reduction may again
be compared with the completion by the same Galois comparison and
Tamagawa criterion as in B.11f. Alternatively all constructions below
can be made over the original henselian trait.

Suppose good marked reduction holds. The finite-cover construction
in B.11d, up to the application of AIK, uses only henselian lifting
and the prime-to-$p$ smooth-pair and Kummer comparisons, so it applies
to the present residue field. Explicitly, a smooth nonempty special
open has a point over a finite separable residue extension: use an
etale chart to affine space over the separable closure, then descend
the point. Henselian lifting and an unramified extension provide an
integral point. If there are at least two cusps, the modulo-$\ell$
abelian boundary quotient is nonzero at every cusp. If there is only
one cusp, positive genus and the Heisenberg $\ell$-group quotient
make its commutator nonzero. A characteristic refinement remains a
finite $\ell$-group quotient. The actual action based at the integral
point is unramified on this quotient; an additional finite unramified
extension kills its finite action and gives a constant-deck Galois
cover with a point above the chosen point. The genus estimate
(B.11d.2), with $\ell\ge5$, makes its compactification have genus
at least two. The regular tame charts $z=w^{e_c}$ give a good
proper model of that compactification. Each of these steps uses the
same SGA~1 statements and proofs specified in B.11d and does not use
residue perfection. B.11f now supplies $\mathcal P_\ell$.

Conversely, B.11f gives good reduction of the proper covering curve
$C$. Its constant deck group $Q$ extends faithfully to the special
fibre of its smooth stable model, by the finite unramified Isom
scheme of \citet[Theorem 1.11]{BC_DM}. Its order is invertible on
the trait. We spell out the one local issue needed over an imperfect
residue field in the tame quotient argument of B.11d.

At a geometric fixed point the stabilizer $H$ embeds faithfully in
the multiplicative group of the tangent line. A prime-to-$p$
finite-order automorphism with derivative one would have its first
nonzero higher coefficient multiplied by its order under iteration;
it must therefore be the identity on the formal local curve, hence
on the connected special curve. Faithfulness excludes this. Thus
$H$ is cyclic; if nontrivial its invariant tangent subspace is zero.
The fixed locus of $H$ is formally smooth over $R_L$ near this point.
Indeed, for a square-zero extension of test rings, smoothness first
gives a lift of a fixed point. The differences of its $H$-translates
form a cocycle in the relative tangent module; averaging, since
$|H|$ is invertible, changes the lift to an $H$-fixed lift. This
establishes the infinitesimal lifting criterion. The zero invariant
tangent space makes the fixed locus etale there. In particular its
residue extension is separable, so the required fixed section is
available after an etale base change even when the original residue
field is imperfect.

At that section the eigenparameter construction of B.11d gives
$h(w)=\zeta w$ and
\[
                       R_L[[w]]^H=R_L[[w^{|H|}]].
\]
The same computation with the semilocal product over an orbit gives
the local ring of the projective quotient $\mathcal C/Q$.
Thus this quotient is smooth and its reduced branch locus consists
etale-locally of disjoint sections. Every special fixed point lifts,
so no branch point occurs only in the special fibre. To check that
different branch sections cannot collide, translate their ramification
points by the deck group so that any alleged collision occurs at the
same special point of $\mathcal C$. Its cyclic stabilizer contains
both nontrivial generic stabilizers, since their fixed equations
persist under specialization. In the eigenparameter $w$, every
nontrivial subgroup of this tame stabilizer has the same unique
local fixed section $w=0$. The two purported distinct ramification
sections must therefore coincide; so do their images. Properness makes
the reduced branch divisor finite etale over $R_L$.
Its generic fibre is exactly $D_L$, since the cover is etale over
$U_L$ and ramified at every cusp. The quotient is therefore a good
marked model of $(\bar U_L,D_L)$.

Finally $L/K$ is unramified, so the identified inertia groups are
equal. Tamagawa's prime-to-$p$ criterion descends good marked
reduction to $R$. If one made the constructions after completion,
the same comparison used at the start descends the conclusion to
the original henselian trait. This proves the criterion in the
full stated scope.
\end{proof}
